\chapter{场景仿真、方案对比与策略验证}

\noindent\fbox{\parbox{0.94\textwidth}{
\textbf{章节状态：待仿真结果回填。}本章当前只冻结测试架构、比较边界和结果表位，不填入
未经重新运行确认的提升比例。所有结果须由V5模型、可追溯输入数据和统一统计脚本生成。}}

\section{数据、场景与验证边界}

\subsection{48 h模型验证场景}

独立48 h短测试仅用于检查多源出力生成、功率守恒、SOC和氢库存状态转移、算力任务约束、
逐时E/H/C分配、故障状态机及因果信息边界，不作为方案经济性比较场景。

\placeholder{48 h验证数据来源、异常注入、守恒误差、状态边界及单元测试结果。}

\subsection{年度主测试场景}

年度场景应覆盖8760 h连续资源、价格、负荷、检修及极端事件。供能端必须由风、光及按
场景启用的海洋能资源序列和设备曲线生成，不预设固定发电量；风电主体占比、容量和设备
可用率均按第3、4章口径设置。每一时段的E/H/C及其他用能比例均实时生效，不得依据事后
得知的年度总供电量回算。

\placeholder{年度数据说明、清洗规则、缺失值、预测误差、容量情景及来源类型。}

\subsection{正常典型48 h采样}

方案对比所用典型48 h从年度正常运行场景中按明确规则抽取，不使用模型开发阶段的48 h
短测试。采样须说明季节、资源水平、价格水平、负荷特征和是否包含外送约束，避免仅选择
对协同方案有利的窗口。

\placeholder{典型窗口选择算法、候选窗口比较及代表性检验。}

\subsection{台风与其他极端事件}

年度场景至少包含台风预警、主动降载/停机、生存状态、关键负荷保障、储能预留、事件后
检查和分级复归。事件参数来自历史气象、设备商或企业应急规则；无来源数值统一标注为
\assumptiontag。

\placeholder{极端事件时序、触发条件、设备状态、恢复策略及动态验证接口。}

\section{单一外送基准与固定策略集合}

先比较单一电力外送基准方案与22种给定固定策略。所有方案使用相同资源、容量、设备可用
率、价格、合同和计量边界；差异仅来自策略规则。固定比例在每个实时小时作用，不提前
获知48 h或全年总电量。

\placeholder{基准规则、22种策略定义、可行性修复顺序和统一输入清单。}

\section{储能、制氢、算力与混合策略对比}

固定策略按“仅储能增强、制氢参与、算力参与、氢算协同、多端口混合”等类别归组，分别
解释资源消纳、产品履约、设备利用和成本变化。模型滚动优化策略不与固定策略混在同一
名次表中，而作为独立的“可部署决策策略”报告；完美预见结果如保留，仅作为理论上界。

\placeholder{固定策略分组表、逐时约束可行率、结果统计及模型策略独立表。}

\section{正常典型48 h运行机理分析}

本节展示年度正常场景中典型48 h的多源可用功率、实际发电、外送、电池充放、制氢、
氢库存、算力服务、海洋负荷和弃能。图表必须同时展示约束触发和策略解释，避免只展示
累计收益。

\placeholder{典型48 h功率堆叠图、SOC/氢库存、E/H/C实时占比、价格与任务状态。}

\section{经济、消纳、可靠与环境绩效}

统一报告单位资源综合价值、全生命周期经济指标、弃能率、产品履约率、算力服务缺口、
EENS/LOLE等适用可靠性指标以及生命周期环境指标。小时级调度结果不得直接宣告构网、
频率、电压或黑启动合格；动态验证结果须单列证据等级。

\placeholder{基准与固定策略结果、模型滚动策略结果、置信区间及数据来源。}

\section{敏感性、边界与稳健性}

敏感性至少覆盖离岸距离与外送成本、可再生资源、设备可用率、储能与电解槽容量、氢氨
和算力合同、预测误差、融资参数及极端事件持续时间。应识别结论翻转点，而非只给出单向
变化曲线。

\placeholder{单因素/多因素敏感性、阈值、策略切换边界及稳健性检验。}

\section{年度因果滚动策略与最优适用区间}

采用48 h预测窗口在全年逐时滚动，每次只执行首个时段决策，再以新实测状态和新预测
更新。策略应包含正常、预测偏差、通信降级、设备故障和台风等模式，并与完美预见上界
区分。最终形成“资源—距离—价格—容量—可靠性”条件下的推荐策略区间，为第6章投资
闸门和第7章结论提供唯一量化输入。

\placeholder{全年滚动结果、极端事件恢复、完美预见差距、适用区间及可复现实验信息。}

